%%%%%%%%%%%%%%%%
%%% num 5 %%%
%%%%%%%%%%%%%%%%

\Chapter{5}

\Verse{1}
{
	\hP{וַיְדַבֵּר יְהֹוָה אֶל מֹשֶׁה לֵּאמֹר׃}
}
{
	\eP{5 And YHWH spoke to Moses, saying, [image]}
}
{
}

\Verse{2}
{
	\hP{צַו אֶת בְּנֵי יִשְׂרָאֵל וִישַׁלְּחוּ מִן הַמַּחֲנֶה כּל צָרוּעַ וְכל זָב וְכֹל טָמֵא לָנָפֶשׁ׃}
}
{
	\eP{
		“Command the children of Israel that they shall have every
		leper and everyone who has an emission and everyone who is
		impure by a person go from the camp. {[image]}
	}
}
{
%	\footnote{
%		impure by a person. Meaning: something that has become impure because of
%		contact with the dead body of a person.
%	}
}

\Verse{3}
{
	\hP{מִזָּכָר עַד נְקֵבָה תְּשַׁלֵּחוּ אֶל מִחוּץ לַמַּחֲנֶה תְּשַׁלְּחוּם וְלֹא יְטַמְּאוּ אֶת מַחֲנֵיהֶם אֲשֶׁר אֲנִי שֹׁכֵן בְּתוֹכָם׃}
}
{
	\eP{
		Male or female, you shall have them go; you shall have
		them go outside the camp so they will not make their
		camps, among which I tent, impure.” {[image]}
	}
}
{
}

\Verse{4}
{
	\hP{וַיַּעֲשׂוּ כֵן בְּנֵי יִשְׂרָאֵל וַיְשַׁלְּחוּ אוֹתָם אֶל מִחוּץ לַמַּחֲנֶה כַּאֲשֶׁר דִּבֶּר יְהֹוָה אֶל מֹשֶׁה כֵּן עָשׂוּ בְּנֵי יִשְׂרָאֵל׃}
}
{
	\eP{
		And the children of Israel did so, and they had them go
		outside the camp. As YHWH spoke to Moses, so the children
		of Israel did. {[image]}
	}
}
{
}

\Verse{5}
{
	\hP{וַיְדַבֵּר יְהֹוָה אֶל מֹשֶׁה לֵּאמֹר׃}
}
{
	\eP{And YHWH spoke to Moses, saying, [image]}
}
{
}

\Verse{6}
{
	\hP{דַּבֵּר אֶל בְּנֵי יִשְׂרָאֵל אִישׁ אוֹ אִשָּׁה כִּי יַעֲשׂוּ מִכּל חַטֹּאת הָאָדָם לִמְעֹל מַעַל בַּיהֹוָה וְאָשְׁמָה הַנֶּפֶשׁ הַהִוא׃}
}
{
	\eP{
		“Speak to the children of Israel: A man or a woman—when
		they will do any of the sins of humans, to make a breach
		against YHWH—so that person has guilt, {[image]}
	}
}
{
}

\Verse{7}
{
	\hP{וְהִתְוַדּוּ אֶת חַטָּאתָם אֲשֶׁר עָשׂוּ וְהֵשִׁיב אֶת אֲשָׁמוֹ בְּרֹאשׁוֹ וַחֲמִישִׁתוֹ יֹסֵף עָלָיו וְנָתַן לַאֲשֶׁר אָשַׁם לוֹ׃}
}
{
	\eP{
		then they shall confess their sin that they have done. And
		he shall pay back for his guilt with its principal, and he
		shall add to it a fifth of it, and he shall give it to the
		one whom he wronged. {[image]}
	}
}
{
}

\Verse{8}
{
	\hP{וְאִם אֵין לָאִישׁ גֹּאֵל לְהָשִׁיב הָאָשָׁם אֵלָיו הָאָשָׁם הַמּוּשָׁב לַיהֹוָה לַכֹּהֵן מִלְּבַד אֵיל הַכִּפֻּרִים אֲשֶׁר יְכַפֶּר בּוֹ עָלָיו׃}
}
{
	\eP{
		And if the man does not have a redeemer, to pay him back
		for the guilt, the guilt that is paid back is YHWH’s, for
		the priest, apart from the ram of atonement by which he
		will make atonement for him. {[image]}
	}
}
{
}

\Verse{9}
{
	\hP{וְכל תְּרוּמָה לְכל קדְשֵׁי בְנֵי יִשְׂרָאֵל אֲשֶׁר יַקְרִיבוּ לַכֹּהֵן לוֹ יִהְיֶה׃}
}
{
	\eP{
		And every donation of all the holy things of the children
		of Israel that they will bring forward to the priest shall
		be his. {[image]}
	}
}
{
}

\Verse{10}
{
	\hP{וְאִישׁ אֶת קֳדָשָׁיו לוֹ יִהְיוּ אִישׁ אֲשֶׁר יִתֵּן לַכֹּהֵן לוֹ יִהְיֶה׃}
}
{
	\eP{
		And each man’s holy things shall be his; what each man
		gives to the priest shall be his.” {[image]}
	}
}
{
%	\footnote{
%		shall be his. Meaning: shall be the priest’s.
%	}
}

\Verse{11}
{
	\hP{וַיְדַבֵּר יְהֹוָה אֶל מֹשֶׁה לֵּאמֹר׃}
}
{
	\eP{And YHWH spoke to Moses, saying, [image]}
}
{
}

\Verse{12}
{
	\hP{דַּבֵּר אֶל בְּנֵי יִשְׂרָאֵל וְאָמַרְתָּ אֲלֵהֶם אִישׁ אִישׁ כִּי תִשְׂטֶה אִשְׁתּוֹ וּמָעֲלָה בוֹ מָעַל׃}
}
{
	\eP{
		“Speak to the children of Israel, and you shall say to
		them: Any man whose wife will go astray and will make a
		breach of faith with him, {[image]}
	}
}
{
%	\footnote{
%		go astray. Hebrew root [image] This situation is known therefore as the
%		case of the [image] Footnote 5:12. go astray. This section deals with
%		the case of the suspected [image]: a woman whose husband suspects her of
%		adultery, but who has not been proved guilty by evidence or witnesses.
%		It would not merit special notice as anything more than any other law,
%		especially since it has not been practiced for at least two millennia
%		(and, I would say, probably more than two-and-a-half millennia), but it
%		has acquired special importance because, first, it is regarded as the
%		only case of trial by ordeal in the Torah and, second, because questions
%		about the legal status of women have risen to prominence in this
%		generation. The woman is made to drink a mixture of holy water, dust
%		from the Tabernacle’s floor, and ink from a document expressing a curse
%		if she is guilty. Then, if “her womb swells and her thigh sags” she is
%		understood to be guilty. There have been many attempts to understand
%		what is going on in this strange case. The traditional Talmudic belief
%		is that a miracle is expected to occur through this procedure. Another
%		view is that the procedure is designed to scare guilty women into
%		confessing rather than face the ceremony and the curse. Another is that,
%		since drinking water with dust and ink does not produce these symptoms,
%		this is a law that is designed to find all women not guilty—and thus
%		protect suspected women from their husbands’ or town’s vengeance.
%		Another view is that this procedure can result in the condition of a
%		prolapsed uterus. All the views leave questions: Why not have some
%		procedure like this for a man who is suspected of adultery? In fact, why
%		not have a procedure like this for any other crime, by a man or a woman?
%		Also, even if the woman is shown guilty, she is not executed, which
%		elsewhere is the penalty for adultery (Lev 20:10; Deut 22:22). So, to
%		begin with: this procedure must relate specifically to something (1)
%		about women and not men, and (2) about sex and adultery. I believe that
%		all of the past explanations are mistaken, and I offer the following
%		explanation of this case: If we were told that a woman’s “womb is
%		swelling and her thigh is sagging,” nearly 100 percent of us would
%		understand this to mean: she is pregnant. So, in this case, it is
%		explicitly about a woman who has had sex with another man in place of
%		her husband. That is, she has not slept with her husband recently, and
%		that is why her pregnancy is proof of adultery. The purpose of drinking
%		the mixture is not to prove her guilt but to bring about the curse
%		through her pregnancy. (Thus it is called “the bitter cursing water”;
%		and the priest says, “Let YHWH make you a curse and an oath among your
%		people when YHWH sets your thigh sagging and your womb swelling.”) This
%		is why the law applies only to women and only to the matter of adultery:
%		it is the only offense in the Torah in which guilt can be determined by
%		this biological fact, even when there are no witnesses and no other
%		evidence. And this is why the woman cannot be executed: She has not been
%		proved guilty in front of judges in a court of law through witnesses.
%		She has been shown to have done this in front of priests at the
%		Tabernacle. Therefore, she “shall bear her crime” (see the comment on
%		5:31). This has a more general implication as well. The reason that we
%		have had such difficulty understanding this case for millennia is that
%		law is constructed on analogy: we decide cases by precedents, by
%		consistent principles that apply to cases in a variety of circumstances.
%		But pregnancy is a condition that is without analogy in human
%		experience. That is why this case reads like no other. And that is why
%		there are agonizing, passionate debates in the present age over whether
%		to save the mother or the baby in a medical crisis during labor, over
%		rights to sperm and eggs that have been preserved, over surrogate
%		mothers, and, above all, over abortion: with arguments about when the
%		developing being inside the womb is a living human in its own right, and
%		about the legal rights of the mother over her own body in balance with
%		the legal protection of the fetus. There simply is no other experience
%		in which two lives are so utterly interconnected. Why are the rabbis of
%		the Talmudic era so unclear about the actual procedure if it had been
%		practiced in recent memory? Because, I think, it had not been practiced
%		since the time of the First Temple in Jerusalem. It required dust from
%		the Tabernacle floor, and the Tabernacle was only in the First Temple,
%		not in the Second. (A report in the Mishnah that Rabbi Yohanan ben
%		Zakkai abolished this procedure makes no sense historically. By the time
%		of his authority, even the Second Temple was gone.) And so it came about
%		that the actual meaning of this ceremony was forgotten after many
%		centuries of not being practiced, and the result has been uncertainty
%		and confusion about its meaning.
%	}
}

\Verse{13}
{
	\hP{וְשָׁכַב אִישׁ אֹתָהּ שִׁכְבַת זֶרַע וְנֶעְלַם מֵעֵינֵי אִישָׁהּ וְנִסְתְּרָה וְהִיא נִטְמָאָה וְעֵד אֵין בָּהּ וְהִוא לֹא נִתְפָּשָׂה׃}
}
{
	\eP{
		and a man has lain with her—an intercourse of seed—and it
		has been hidden from her husband’s eyes, and she has kept
		concealed, and she has been made impure, and there is no
		witness against her, and she has not been caught, {[image]}
	}
}
{
}

\Verse{14}
{
	\hP{וְעָבַר עָלָיו רוּחַ קִנְאָה וְקִנֵּא אֶת אִשְׁתּוֹ וְהִוא נִטְמָאָה אוֹ עָבַר עָלָיו רוּחַ קִנְאָה וְקִנֵּא אֶת אִשְׁתּוֹ וְהִיא לֹא נִטְמָאָה׃}
}
{
	\eP{
		and a spirit of jealousy has come over him, and he is
		jealous about his wife, and she has been made impure, or a
		spirit of jealousy has come over him, and he is jealous
		about his wife, and she has not been made impure: {[image]}
	}
}
{
}

\Verse{15}
{
	\hP{וְהֵבִיא הָאִישׁ אֶת אִשְׁתּוֹ אֶל הַכֹּהֵן וְהֵבִיא אֶת קרְבָּנָהּ עָלֶיהָ עֲשִׂירִת הָאֵיפָה קֶמַח שְׂעֹרִים לֹא יִצֹק עָלָיו שֶׁמֶן וְלֹא יִתֵּן עָלָיו לְבֹנָה כִּי מִנְחַת קְנָאֹת הוּא מִנְחַת זִכָּרוֹן מַזְכֶּרֶת עָוֺן׃}
}
{
	\eP{
		then the man shall bring his wife to the priest and shall
		bring her offering along with her: a tenth of an ephah of
		barley flour. He shall not pour oil on it and shall not
		put frankincense on it, because it is an offering about
		jealousies, an offering of bringing to mind: causing a
		crime to be brought to mind. {[image]}
	}
}
{
}

\Verse{16}
{
	\hP{וְהִקְרִיב אֹתָהּ הַכֹּהֵן וְהֶעֱמִדָהּ לִפְנֵי יְהֹוָה׃}
}
{
	\eP{
		“And the priest shall bring her forward and stand her in
		front of YHWH. {[image]}
	}
}
{
}

\Verse{17}
{
	\hP{וְלָקַח הַכֹּהֵן מַיִם קְדֹשִׁים בִּכְלִי חָרֶשׂ וּמִן הֶעָפָר אֲשֶׁר יִהְיֶה בְּקַרְקַע הַמִּשְׁכָּן יִקַּח הַכֹּהֵן וְנָתַן אֶל הַמָּיִם׃}
}
{
	\eP{
		And the priest shall take holy water in a clay container,
		and the priest shall take some of the dust that will be on
		the Tabernacle’s floor and put it into the water. {[image]}
	}
}
{
%	\footnote{
%		holy water. This is the only reference to holy water in the Torah. It
%		may be the water from the basin that the priests use for washing (Exod
%		30:17–21); the basin is among the things that are identified as “holy of
%		holies” (30:28–29). Footnote 5:17. dust. The woman must drink the water
%		containing this dust (5:24). Some have claimed that there is a parallel
%		between this case and the episode of the golden calf. Moses grinds the
%		golden calf thin as dust (Exod 32:20; Deut 9:21), scatters it in water,
%		and makes the people drink (Exodus 32). My reaction is that such a view
%		is tenuous at best. The people in Exodus are guilty! The [image] case is
%		a process designed to determine guilt. And the dust from the Tabernacle
%		is something positive, while the dust of the calf is the opposite. It is
%		true that the people of Israel’s turning to pagan gods is sometimes
%		compared to adultery in the Tanak: Israel is described as a faithless
%		wife. And the term that means “a breach of faith” (ma‘al) with a husband
%		here refers to breaches of faith with God in all its other occurrences
%		in the Torah. But the cases of the [image] and the golden calf have so
%		many fundamental differences that drawing a connection between them may
%		mislead more than shed light.
%	}
}

\Verse{18}
{
	\hP{וְהֶעֱמִיד הַכֹּהֵן אֶת הָאִשָּׁה לִפְנֵי יְהֹוָה וּפָרַע אֶת רֹאשׁ הָאִשָּׁה וְנָתַן עַל כַּפֶּיהָ אֵת מִנְחַת הַזִּכָּרוֹן מִנְחַת קְנָאֹת הִוא וּבְיַד הַכֹּהֵן יִהְיוּ מֵי הַמָּרִים הַמְאָרְרִים׃}
}
{
	\eP{
		And the priest shall stand the woman in front of YHWH and
		loosen the hair of the woman’s head and put the offering
		of bringing to mind on her hands; it is an offering of
		jealousies. And in the priest’s hand shall be the bitter
		cursing water. {[image]}
	}
}
{
}

\Verse{19}
{
	\hP{וְהִשְׁבִּיעַ אֹתָהּ הַכֹּהֵן וְאָמַר אֶל הָאִשָּׁה אִם לֹא שָׁכַב אִישׁ אֹתָךְ וְאִם לֹא שָׂטִית טֻמְאָה תַּחַת אִישֵׁךְ הִנָּקִי מִמֵּי הַמָּרִים הַמְאָרְרִים הָאֵלֶּה׃}
}
{
	\eP{
		And the priest shall have her swear, and he shall say to
		the woman: ‘If a man has not lain with you, and if you
		have not gone astray in impurity with someone in place of
		your husband, be cleared by this bitter cursing water.
		{[image]}
	}
}
{
%	\footnote{
%		in place of your husband. Suggesting that if she is pregnant it must be
%		by the other man, not the husband. Those who have other views of the
%		[image] procedure understand this phrase to mean “under your husband,”
%		and they interpret it as meaning “under his authority.” But there is no
%		reason why that had to be mentioned in the context of these
%		instructions. Moreover, the term here, Hebrew [image], means “in place
%		of” sixteen times in the Torah but never means “under someone’s
%		authority” (the term for that is “under the hand of,” Hebrew [image]
%		yad; cf. Gen 41:35).
%	}
}

\Verse{20}
{
	\hP{וְאַתְּ כִּי שָׂטִית תַּחַת אִישֵׁךְ וְכִי נִטְמֵאת וַיִּתֵּן אִישׁ בָּךְ אֶת שְׁכבְתּוֹ מִבַּלְעֲדֵי אִישֵׁךְ׃}
}
{
	\eP{
		But you, if you have gone astray with someone in place of
		your husband and if you have been made impure, and if a
		man other than your husband has had his intercourse with
		you,’ {[image]}
	}
}
{
}

\Verse{21}
{
	\hP{וְהִשְׁבִּיעַ הַכֹּהֵן אֶת הָאִשָּׁה בִּשְׁבֻעַת הָאָלָה וְאָמַר הַכֹּהֵן לָאִשָּׁה יִתֵּן יְהֹוָה אוֹתָךְ לְאָלָה וְלִשְׁבֻעָה בְּתוֹךְ עַמֵּךְ בְּתֵת יְהֹוָה אֶת יְרֵכֵךְ נֹפֶלֶת וְאֶת בִּטְנֵךְ צָבָה׃}
}
{
	\eP{
		then the priest shall have the woman swear a curse oath,
		and the priest shall say to the woman, ‘let YHWH make you
		a curse and an oath among your people when YHWH sets your
		thigh sagging and your womb swelling, {[image]}
	}
}
{
%	\footnote{
%		womb. Hebrew been. This word is commonly understood to mean the woman’s
%		“belly,” but in every one of its other ten occurrences in the Torah it
%		means “womb” (see, e.g. Gen 25:23). This is further evidence that
%		pregnancy is what this is all about.
%	}
}

\Verse{22}
{
	\hP{וּבָאוּ הַמַּיִם הַמְאָרְרִים הָאֵלֶּה בְּמֵעַיִךְ לַצְבּוֹת בֶּטֶן וְלַנְפִּל יָרֵךְ וְאָמְרָה הָאִשָּׁה אָמֵן אָמֵן׃}
}
{
	\eP{
		and this cursing water will come in your insides, to swell
		the womb and make the thigh sag.’ “And the woman shall
		say, ‘Amen, amen.’ {[image]}
	}
}
{
%	\footnote{
%		to swell the womb and make the thigh sag. It is not clear whether this
%		refers to God’s making this happen (as is explicit in the preceding
%		verse) or to the water’s making it happen. It appears to me that God is
%		understood here to bring about the pregnancy—in the sense that all
%		pregnancies are understood to be enabled by God—and that the water
%		brings about the curse: “the bitter cursing water.”
%	}
}

\Verse{23}
{
	\hP{וְכָתַב אֶת הָאָלֹת הָאֵלֶּה הַכֹּהֵן בַּסֵּפֶר וּמָחָה אֶל מֵי הַמָּרִים׃}
}
{
	\eP{
		“And the priest shall write these curses in a scroll and
		rub them into the bitter water. {[image]}
	}
}
{
}

\Verse{24}
{
	\hP{וְהִשְׁקָה אֶת הָאִשָּׁה אֶת מֵי הַמָּרִים הַמְאָרְרִים וּבָאוּ בָהּ הַמַּיִם הַמְאָרְרִים לְמָרִים׃}
}
{
	\eP{
		And he shall have the woman drink the bitter cursing
		water, and the bitter cursing water will come into her.
		{[image]}
	}
}
{
%	\footnote{
%		the bitter cursing water. Hebrew mê hammrîm hamme’rrîm. The obvious
%		alliteration is one of several cases of wordplay in this text: Hebrew
%		[image] (it has been hidden) reverses the root letters of [image]
%		(breach). Hebrew [image] (a tenth) reverses the root letters of [image]
%		(barley). Hebrew [image] (loosen the hair) reverses [image] (dust). And
%		Hebrew [image] (cursing) has strong alliteration with [image] (bitter).
%		This wordplay involves much more than artistry with language. Each of
%		these puns is linked to an aspect of this ceremony that is unique in
%		some way: This is the only case in which [image] (to make a breach) is
%		used for an offense against a human rather than God (see the comment on
%		Lev 5:15). It is the only time that [image] (barley flour) is used for a
%		grain offering. It is the only use of [image] the It is the only time
%		that (dust) from the Tabernacle. And this strange use of water (holy
%		water, with a curse melted into it) is the most unusual of all. Puns are
%		common in the Torah’s stories, but their appearance in a law is
%		extremely unusual and unexpected. They highlight the unusual aspect of
%		this case, and they summon to mind associations with the realm of magic,
%		in which the artistry of language and poetry figure in incantations and
%		ceremonies. Likewise, the dissolving of the words of the curse into the
%		water, the dust from the Tabernacle, the water from the priests’ basin,
%		the act of drinking a mixture that produces a curse: much of this has
%		parallels in magic in other cultures, and all of this renders this
%		ceremony mysterious. And it fits with the fact that the function of this
%		procedure is not to establish legal guilt but rather to produce the
%		curse.
%	}
}

\Verse{25}
{
	\hP{וְלָקַח הַכֹּהֵן מִיַּד הָאִשָּׁה אֵת מִנְחַת הַקְּנָאֹת וְהֵנִיף אֶת הַמִּנְחָה לִפְנֵי יְהֹוָה וְהִקְרִיב אֹתָהּ אֶל הַמִּזְבֵּחַ׃}
}
{
	\eP{
		And the priest shall take the offering of jealousies from
		the woman’s hand and elevate the offering in front of
		YHWH, and he shall bring her forward to the altar. {[image]}
	}
}
{
}

\Verse{26}
{
	\hP{וְקָמַץ הַכֹּהֵן מִן הַמִּנְחָה אֶת אַזְכָּרָתָהּ וְהִקְטִיר הַמִּזְבֵּחָה וְאַחַר יַשְׁקֶה אֶת הָאִשָּׁה אֶת הַמָּיִם׃}
}
{
	\eP{
		And the priest shall take a fistful from the grain
		offering, a representative portion of it, and burn it to
		smoke at the altar, and after that he shall have the woman
		drink the water. {[image]}
	}
}
{
}

\Verse{27}
{
	\hP{וְהִשְׁקָהּ אֶת הַמַּיִם וְהָיְתָה אִם נִטְמְאָה וַתִּמְעֹל מַעַל בְּאִישָׁהּ וּבָאוּ בָהּ הַמַּיִם הַמְאָרְרִים לְמָרִים וְצָבְתָה בִטְנָהּ וְנָפְלָה יְרֵכָהּ וְהָיְתָה הָאִשָּׁה לְאָלָה בְּקֶרֶב עַמָּהּ׃}
}
{
	\eP{
		When he has had her drink the water, then it will be, if
		she has been made impure and has made a breach of faith
		with her husband, when the bitter cursing water will come
		into her, and her womb will swell and her thigh will sag,
		then the woman will become a curse among her people.
		{[image]}
	}
}
{
}

\Verse{28}
{
	\hP{וְאִם לֹא נִטְמְאָה הָאִשָּׁה וּטְהֹרָה הִוא וְנִקְּתָה וְנִזְרְעָה זָרַע׃}
}
{
	\eP{
		And if the woman has not been made impure, and she is
		pure, then she shall be cleared and shall conceive seed.
		{[image]}
	}
}
{
}

\Verse{29}
{
	\hP{זֹאת תּוֹרַת הַקְּנָאֹת אֲשֶׁר תִּשְׂטֶה אִשָּׁה תַּחַת אִישָׁהּ וְנִטְמָאָה׃}
}
{
	\eP{
		“This is the instruction for jealousies, when a woman will
		go astray with someone in place of her husband and be made
		impure, {[image]}
	}
}
{
}

\Verse{30}
{
	\hP{אוֹ אִישׁ אֲשֶׁר תַּעֲבֹר עָלָיו רוּחַ קִנְאָה וְקִנֵּא אֶת אִשְׁתּוֹ וְהֶעֱמִיד אֶת הָאִשָּׁה לִפְנֵי יְהֹוָה וְעָשָׂה לָהּ הַכֹּהֵן אֵת כּל הַתּוֹרָה הַזֹּאת׃}
}
{
	\eP{
		or when a spirit of jealousy will come over a man and he
		will be jealous about his wife. And the priest shall stand
		the woman in front of YHWH and do all of this instruction
		to her. {[image]}
	}
}
{
}

\Verse{31}
{
	\hP{וְנִקָּה הָאִישׁ מֵעָוֺן וְהָאִשָּׁה הַהִוא תִּשָּׂא אֶת עֲוֺנָהּ׃}
}
{
	\eP{
		And the man shall be clear of a crime, and that woman
		shall bear her crime.” {[image]}
	}
}
{
%	\footnote{
%		the man shall be clear of a crime. This is usually understood to mean
%		that the husband is innocent of the crime of making a false charge
%		against his wife. But this is a questionable reading since the text does
%		not say that he has made any charge against her; he has just brought her
%		to the priest to make the determination of whether his jealousy is
%		correct or not. This may rather refer to the other man, the one who had
%		intercourse with the woman. The procedure proves only that the woman has
%		had intercourse; it cannot prove who the man was. It is therefore not
%		possible to convict the man by this procedure. This also answers the
%		question of why this procedure applies only to women. A man accused of
%		adultery does not go through this. Because only women can become
%		pregnant, it is only women who can be proved to have committed adultery
%		in the absence of witnesses or any other evidence. (This brings to mind
%		the case of Bathsheba, who becomes pregnant through adultery with David.
%		Her pregnancy while her husband has been away at war is what would make
%		her adultery known; 2 Samuel 11.)
%	}
}
